% ===========================================================================
\section{Conclusions}\label{sec:conclusions}
% ===========================================================================

We have written a modular Einstein--Boltzmann solver and used it to compute the CMB temperature and polarization
spectra and the matter power spectrum for a flat $\Lambda$CDM model. The four milestones gave:
\begin{enumerate}
\item The background. Radiation--matter equality is at $z=3400$ and acceleration starts at $z=0.63$, before
      matter--$\Lambda$ equality at $z=0.29$, because the negative pressure of $\Lambda$ counts twice in the
      acceleration equation. The Universe is $13.85$\,Gyr old, older than a matter-only universe could be. Supernova
      distances require an accelerating expansion ($\chi^2=31$ versus 317 for a matter-only universe) and prefer
      $h\approx0.70$ when the density parameters are held fixed.
\item The recombination history. The huge number of photons per baryon delays recombination until
      $k_BT\approx\epsilon_0/50$. Photons decouple at $z=1090$, when the Universe is $372\,000$ years old (or at
      $z_*=1085$ including reionization). The baryons, outnumbered by the photons, follow at $z_d=1060$, with sound
      horizon $147\Mpc$. Saha equilibrium ignores the $n=2$ bottleneck and would predict decoupling $\sim200$ units of
      redshift too early and no freeze-out.
\item The perturbations show the main physics of structure formation: acoustic oscillations of the photon--baryon
      fluid, driven by decaying potentials and damped by photon diffusion; the stagnation of dark matter growth in the
      radiation era; baryons falling into dark matter wells after decoupling; the decay of the potentials in the
      $\Lambda$ era; and the brief generation of polarization during decoupling.
\item The spectra reproduce the Planck TT, EE and TE data and the SDSS galaxy power spectrum without any fitting. The
      CMB peaks are the acoustic oscillations frozen at decoupling and projected over $\chi_*$, with spacing set by
      $\ell_A=\pi\chi_*/r_s=303$. Their heights encode the baryon and matter densities, and the turnover of $P(k)$ at
      $k_{\rm eq}$ is the memory of radiation--matter equality. The oscillating residuals of the unlensed TT spectrum are
      explained by gravitational lensing: with our lensing module the $\chi^2$ against the 83 Planck high-$\ell$ bins
      falls from 781 to 220, and the lensed spectra and the lensing B modes agree with CLASS.
\end{enumerate}
All modules agree with CLASS at the percent level or better. As an extension, the code also handles open and
closed universes (Sec.~\ref{sec:curvature}), with hyperspherical Bessel functions and the discrete modes of a closed
space; for $-0.2\le\Om_K\le0.2$ it agrees with CLASS as well as in the flat case, and the first acoustic peak
follows the geometric prediction $\ell_{\rm peak}\simeq0.73\,\pi r(\chi_*)/r_s$. With a massive neutrino species
(Sec.~\ref{sec:neutrinos}) the code reproduces the free-streaming suppression of the matter power spectrum,
$\approx-7f_\nu$ on small scales, and of the lensing potential, with the effect of the mass agreeing with CLASS to
$5\times10^{-4}$ in $P(k)$ and $C_\ell^{TT}$.
The two-point functions of Sec.~\ref{sec:twopoint} show the acoustic physics in real space: a shoulder near $1^\circ$
in the CMB temperature correlation, alternating radial and tangential polarization around hot spots, and the BAO
shell at $r_d$ in the matter correlation function and, projected, at $r_d/\chi$ in $w(\theta)$.

The recombination model was the largest remaining physical approximation. With RECFAST (Sec.~\ref{sec:recfast}),
$X_e$ agrees with CLASS's RECFAST to $3.3\times10^{-4}$, and the TT and EE spectra agree to a few
$10^{-3}$ (rms) up to $\ell=2000$, compared with 1--3\% for the Peebles model. The model library of
Sec.~\ref{sec:models} extends the code to eleven dark-energy and modified-gravity models. Each is validated against
its $\Lambda$CDM limit, the unused constraint equations and independent calculations, and against CLASS, hi\_class
and CLASS\_EDE where these apply. Natural next steps are:
\begin{itemize}
\item Nonlinear corrections to $P(k)$ (halofit or HMcode) for weak-lensing and galaxy-clustering predictions, and
      their extension to modified gravity.
\item Horndeski gravity with $\alpha_T\neq0$, and models built directly from the Horndeski functions $G_i(\phi,X)$.
\item Several massive neutrino species with different masses (the code has one), and $B$-mode polarization from
      tensor modes.
\item MCMC parameter estimation with the full code; the MCMC fits of the course project are omitted from this report
      as instructed.
\end{itemize}
The code takes about one minute to run on a laptop, which is fast enough for such exploratory work.

\begin{acknowledgments}
We thank the authors of CLASS for making their code public, which was invaluable for validation. This work used the
GNU Scientific Library, NumPy, SciPy and Matplotlib.
\end{acknowledgments}
